% =============================================================
%  Section 7 – Robustness
%  Label: sec:robustness
% =============================================================

\section{Robustness}\label{sec:robustness}

This section examines inference with a small number of clusters,
aggregate sign replication, sensitivity to fixed-effect structure,
and the 2016 CPI weight revision.

% -------------------------------------------------------------
\subsection{Wild Cluster Bootstrap}\label{sec:wcb}
% -------------------------------------------------------------

The preferred specifications cluster standard errors at the
province level.
China has 31 mainland provinces, municipalities, and autonomous
regions, so the maximum number of clusters available for
inference is around 25 to 31 depending on data coverage.
This cluster count is in the range where cluster-robust
standard errors based on the sandwich formula are known to
over-reject: the asymptotic approximation performs poorly
below approximately 30 to 50 clusters, leading to
$t$-statistics that are too large in finite samples
\citep{Cameron2008, MacKinnon2017}.

To address this, I report wild cluster bootstrap $p$-values
for the main coefficient estimates.
I implement the Webb six-point weight distribution, which
performs well with cluster counts below 30.
All bootstraps use 9,999 replications and impose the null
hypothesis on each coefficient separately.

For the expectation equation under province-plus-wave FE,
the bootstrap $p$-value on $\hat{\beta}_1$ is $0.088$,
compared with the asymptotic $p = 0.100$.
Under region-by-wave FE, the bootstrap $p$-value is $0.237$
(asymptotic: $0.174$).
The expectation coefficient is thus marginal under both
inference methods.

The forecast-error coefficient $\hat{\beta}_3$ is robust
to the bootstrap correction.
Under province-plus-wave FE, bootstrap $p < 0.001$
(asymptotic $p < 0.001$).
Under region-by-wave FE, bootstrap $p < 0.001$
(asymptotic $p < 0.001$).
The core overreaction finding in the forecast-error equation
survives the most conservative available inference with
31 clusters.

% -------------------------------------------------------------
\subsection{Aggregate Sign Check: PBoC Depositor Survey}%
\label{sec:pboc_sign}
% -------------------------------------------------------------

The PBoC Urban Depositor Survey provides an independent
aggregate-level check on the sign of the overreaction pattern.
Using the Carlson--Parkin quantified expectations from the
quarterly depositor survey, I estimate the baseline
revision-error regression at the aggregate level.
The sign of the revision-error slope is negative,
consistent with the household-level overreaction finding in
Section~\ref{sec:results}.

Three caveats govern the interpretation.
First, $N \approx 32$ is too small for precise inference.
Second, the Carlson--Parkin quantification introduces
measurement error.
Third, aggregate time-series regressions cannot rule out
omitted macro confounders.
The aggregate evidence is corroborative only; the
household-level shock design in Section~\ref{sec:results}
is the primary source of identification.

% -------------------------------------------------------------
\subsection{Alternative Fixed Effects}\label{sec:altfe}
% -------------------------------------------------------------

The preferred specifications use region-by-wave fixed
effects, grouping the 31 provincial units into macro
regions and interacting the region indicators with
wave dummies.
This approach preserves within-region variation in shock
intensity while absorbing region-specific macro trends.
I evaluate sensitivity to this choice by estimating
alternative fixed-effect structures.

The first alternative uses province fixed effects and
wave fixed effects entered separately (two-way rather than
interacted).
This specification permits province-specific time-invariant
differences and wave-specific national shocks to be absorbed
independently but does not absorb time-varying
province-specific trends.
The province-and-wave specification yields $\hat{\beta}_1
= 0.090$ ($p = 0.10$) and $\hat{\beta}_3 = -5.62$
($p < 0.001$), close to the preferred region-by-wave
estimates.
The stability across these two FE structures indicates that
the results are not driven by the specific geographic
grouping used for region-by-wave absorption.

% -------------------------------------------------------------
\subsection{CPI Weight Revision}\label{sec:cpi_weights}
% -------------------------------------------------------------

The National Bureau of Statistics revised the CPI basket
weights in 2016.
The revision reduced the headline food weight from
approximately 31 percent to 28 percent and adjusted
sub-weights within the food category.
If these revisions materially shifted the relationship
between meat prices and headline CPI, the pass-through
regression in Section~\ref{sec:eq2} could behave
differently before and after 2016, and the implied forecast
error could reflect a weight-regime change rather than
household overreaction.

Two observations mitigate this concern.
First, the meat CPI data used for shock construction begin in
January 2016, so the estimation sample (CFPS waves 2016--2022)
falls entirely within the post-revision weight regime.
The pre-2016 waves (2012, 2014) lack province-level meat CPI
data and do not enter the main specifications.
There is therefore no weight-regime boundary within the
estimation sample.

Second, the consistency of results across the 2016, 2018,
2020, and 2022 waves, as shown in the wave-specific
meat-shock summary statistics, indicates that the relationship
between meat shocks and expectations is stable within the
post-revision period.

Taken together, the robustness exercises confirm that the
forecast-error result is not an artifact of inference method,
fixed-effect structure, or CPI measurement regime.
The expectation coefficient is less precisely estimated,
which is expected given the demanding fixed-effect structure,
but the core finding that meat-price shocks generate
systematic overprediction of future inflation is robust.
